\documentclass[journal]{vgtc}                % final (journal style)
%\documentclass[review,journal]{vgtc}         % review (journal style)
%\documentclass[widereview]{vgtc}             % wide-spaced review
%\documentclass[preprint,journal]{vgtc}       % preprint (journal style)

%% Uncomment one of the lines above depending on where your paper is
%% in the conference process. ``review'' and ``widereview'' are for review
%% submission, ``preprint'' is for pre-publication, and the final version
%% doesn't use a specific qualifier.

%% Please use one of the ``review'' options in combination with the
%% assigned online id (see below) ONLY if your paper uses a double blind
%% review process. Some conferences, like IEEE Vis and InfoVis, have NOT
%% in the past.

%% Please use the ``preprint''  option when producing a preprint version
%% for sharing your article on an open access repository

%% Please note that the use of figures other than the optional teaser is not permitted on the first page
%% of the journal version.  Figures should begin on the second page and be
%% in CMYK or Grey scale format, otherwise, colour shifting may occur
%% during the printing process.  Papers submitted with figures other than the optional teaser on the
%% first page will be refused. Also, the teaser figure should only have the
%% width of the abstract as the template enforces it.

%% These few lines make a distinction between latex and pdflatex calls and they
%% bring in essential packages for graphics and font handling.
%% Note that due to the \DeclareGraphicsExtensions{} call it is no longer necessary
%% to provide the the path and extension of a graphics file:
%% \includegraphics{diamondrule} is completely sufficient.
%%
\ifpdf%                                % if we use pdflatex
  \pdfoutput=1\relax                   % create PDFs from pdfLaTeX
  \pdfcompresslevel=9                  % PDF Compression
  \pdfoptionpdfminorversion=7          % create PDF 1.7
  \ExecuteOptions{pdftex}
  \usepackage{graphicx}                % allow us to embed graphics files
  \DeclareGraphicsExtensions{.pdf,.png,.jpg,.jpeg} % for pdflatex we expect .pdf, .png, or .jpg files
\else%                                 % else we use pure latex
  \ExecuteOptions{dvips}
  \usepackage{graphicx}                % allow us to embed graphics files
  \DeclareGraphicsExtensions{.eps}     % for pure latex we expect eps files
\fi%

%% it is recomended to use ``\autoref{sec:bla}'' instead of ``Fig.~\ref{sec:bla}''
\graphicspath{{figures/}{pictures/}{images/}{./}} % where to search for the images

\usepackage{microtype}                 % use micro-typography (slightly more compact, better to read)
\PassOptionsToPackage{warn}{textcomp}  % to address font issues with \textrightarrow
\usepackage{textcomp}                  % use better special symbols
\usepackage{mathptmx}                  % use matching math font
\usepackage{times}                     % we use Times as the main font
\renewcommand*\ttdefault{txtt}         % a nicer typewriter font
\usepackage{cite}                      % needed to automatically sort the references
\usepackage{tabu}                      % only used for the table example
\usepackage{booktabs}                  % only used for the table example
%% We encourage the use of mathptmx for consistent usage of times font
%% throughout the proceedings. However, if you encounter conflicts
%% with other math-related packages, you may want to disable it.

%% In preprint mode you may define your own headline. If not, the default IEEE copyright message will appear in preprint mode.
%\preprinttext{To appear in IEEE Transactions on Visualization and Computer Graphics.}

%% In preprint mode, this adds a link to the version of the paper on IEEEXplore
%% Uncomment this line when you produce a preprint version of the article 
%% after the article receives a DOI for the paper from IEEE
%\ieeedoi{xx.xxxx/TVCG.201x.xxxxxxx}

%% If you are submitting a paper to a conference for review with a double
%% blind reviewing process, please replace the value ``0'' below with your
%% OnlineID. Otherwise, you may safely leave it at ``0''.
\onlineid{0}

%% declare the category of your paper, only shown in review mode
\vgtccategory{Research}
%% please declare the paper type of your paper to help reviewers, only shown in review mode
%% choices:
%% * algorithm/technique
%% * application/design study
%% * evaluation
%% * system
%% * theory/model
\vgtcpapertype{please specify}

%% Paper title.
\title{CS594 - Big Data Visualization \& Analytics: Proposal}

%% This is how authors are specified in the journal style

%% indicate IEEE Member or Student Member in form indicated below
\author{Roy G. Biv, Ed Grimley, \textit{Member, IEEE}, and Martha Stewart}
\authorfooter{
%% insert punctuation at end of each item
\item
 Roy G. Biv is with Starbucks Research. E-mail: roy.g.biv@aol.com.
\item
 Ed Grimley is with Grimley Widgets, Inc.. E-mail: ed.grimley@aol.com.
\item
 Martha Stewart is with Martha Stewart Enterprises at Microsoft
 Research. E-mail: martha.stewart@marthastewart.com.
}

%other entries to be set up for journal
\shortauthortitle{Biv \MakeLowercase{\textit{et al.}}: Global Illumination for Fun and Profit}
%\shortauthortitle{Firstauthor \MakeLowercase{\textit{et al.}}: Paper Title}

%% Abstract section.
%\abstract{Duis autem vel eum iriure dolor in hendrerit in vulputate
%velit esse molestie consequat, vel illum dolore eu feugiat nulla
%facilisis at vero eros et accumsan et iusto odio dignissim qui blandit
%praesent luptatum zzril delenit augue duis dolore te feugait nulla
%facilisi. Lorem ipsum dolor sit amet, consectetuer adipiscing elit,
%sed diam nonummy nibh euismod tincidunt ut laoreet dolore magna
%aliquam erat volutpat. Ut wisi enim ad minim veniam, quis nostrud exerci tation ullamcorper
%suscipit lobortis nisl ut aliquip ex ea commodo consequat. Duis autem
%vel eum iriure dolor in hendrerit in vulputate velit esse molestie
%consequat, vel illum dolore eu feugiat nulla facilisis at vero eros et
%accumsan et iusto odio dignissim qui blandit praesent luptatum zzril
%delenit augue duis dolore te feugait nulla facilisi.%
%} % end of abstract

%% Keywords that describe your work. Will show as 'Index Terms' in journal
%% please capitalize first letter and insert punctuation after last keyword
%% \keywords{Radiosity, global illumination, constant time}

%% ACM Computing Classification System (CCS). 
%% See <http://www.acm.org/class/1998/> for details.
%% The ``\CCScat'' command takes four arguments.

%% \CCScatlist{ % not used in journal version
%% \CCScat{K.6.1}{Management of Computing and Information Systems}%
%% {Project and People Management}{Life Cycle};
%% \CCScat{K.7.m}{The Computing Profession}{Miscellaneous}{Ethics}
%% }

%% A teaser figure can be included as follows
%\teaser{
%  \centering
%  \includegraphics[width=\linewidth]{CypressView}
%  \caption{In the Clouds: Vancouver from Cypress Mountain. Note that the teaser may not be wider than the abstract block.}
%  \label{fig:teaser}
%}

%% Uncomment below to disable the manuscript note
\renewcommand{\manuscriptnotetxt}{}

%% Copyright space is enabled by default as required by guidelines.
%% It is disabled by the 'review' option or via the following command:
% \nocopyrightspace


\vgtcinsertpkg

%%%%%%%%%%%%%%%%%%%%%%%%%%%%%%%%%%%%%%%%%%%%%%%%%%%%%%%%%%%%%%%%
%%%%%%%%%%%%%%%%%%%%%% START OF THE PAPER %%%%%%%%%%%%%%%%%%%%%%
%%%%%%%%%%%%%%%%%%%%%%%%%%%%%%%%%%%%%%%%%%%%%%%%%%%%%%%%%%%%%%%%%

\begin{document}

%% The ``\maketitle'' command must be the first command after the
%% ``\begin{document}'' command. It prepares and prints the title block.

%% the only exception to this rule is the \firstsection command
\firstsection{Introduction}

\maketitle


\section{Introduction}
Navigating large code bases is an essential skill for software developers to have.
There are many tools that help developers navigate code bases such as code editors like CLion and PyCharm \cite{jetbrains}, and Google's Code Search \cite{googlecodesearch}. There are different aspects of code bases that a developer may be interested in and there are various tools available for each of them. This project proposes a visualization tool to navigate code bases using a tree visualization of the directory structure. The interactive tree visualization can be further used by the programmer to query information on the files and directories. This project focuses on the version control aspect of code bases.
\section{Related work}
\section{Data description}

The data for this project lies within a cloned git repository. Git stores data as a series of snapshots instead of storing differences between file versions like many other version control systems. Each commit in Git represents a snapshot of the entire working directory at a given point in time. The data is stored as objects, which are identified by a unique SHA-1 hash. There are four types of objects:
\begin{itemize}
    \item Blob: Represents a file and its contents.
    \item Tree: Represents a directory and stores references to blob objects (files) and other tree objects (subdirectories).
    \item Commit: Represents a snapshot of the repository, including metadata such as author, committer, date, and a message. A commit object also references the tree object representing the root directory of the snapshot.
    \item Tag: A lightweight object that references a commit object to create a more human-readable name for a particular commit, typically used for marking releases.
\end{itemize}
Git objects are stored in the ".git/objects" directory as individual files. However, to save space and improve performance, Git periodically repacks objects into packfiles. A packfile contains multiple objects and their delta-compressed versions, which helps reduce storage space.
Although we could write our own scripts to go through all this data stored in the .git folder, it is not necessary to do so. Git already provides a efficient command line interface to query this database. Additional details about the inner workings of git can be found at \cite{gitinternals} and the git docs \cite{gitdocs}. 

Some data does not reside locally in the cloned repository and can be accessed using Github's REST API \cite{gitrestapi}.

A general list of attributes that one can query using git/Github's REST API:
\begin{itemize}
    \item branches
    \item commits
    \item contributors
    \item pull requests
    \item forks
    \item files
    \item directories
    \item code and line count - additions/deletions
\end{itemize}

This list is not perfect as each attribute on the list can have further details such as each commit has a date associated with it and each pull request only affects the changed files and directories. The goal of this project is to use a combination of these attributes and make it easy for a developer to gain insights about the code base. For example, a developer might be interested in the pull requests that have modified a specific file.
\section{Proposal}
\section{Results}
Over the course of this semester I was able to implement the font and back end for rendering the graph and loading different git repositories. This includes the methods discussed in the research section \ref{sec:research}. The project overall gave me an idea on how to handle large graph data and experiment with different UI setups to make it easy for the client to explore the graph.

The next step would be implement the front end for running the queries.
\section{Conclusion}


%\bibliographystyle{abbrv}
\bibliographystyle{abbrv-doi}
%\bibliographystyle{abbrv-doi-narrow}
%\bibliographystyle{abbrv-doi-hyperref}
%\bibliographystyle{abbrv-doi-hyperref-narrow}

\bibliography{proposal}
\end{document}

